机制面板包含六个CAL五复访序列及其嵌套前缀，覆盖两个物理参考场景。
五种固定模型／干预配置共150次独立进程推理，另有2次仅反向传播的梯度诊断。
全部152项已完成，保存的预测与真值已用于独立重算官方状态。
75份汇总状态分别覆盖25个总体和50个场景--长度组合。

训练后的反馈模块在两个梯度诊断中的梯度范数分别为0.9788与1.4578；
交叉注意力、自注意力、解码FFN和mask head的梯度均有限且非零。
诊断使用原始任务目标与12个辅助输出，不进行优化器更新，编码器没有梯度，所有参数保持不变。
这证明模块连接到实际任务目标，不代表其提供准确率收益。
在28组严格上游对齐的配对中，辅助阶段0--7保持完全一致；
反馈后的四个辅助mask输出及最终mask均改变。
紧接反馈的阶段8类别logits不变，阶段9--11及最终类别logits随后改变，
与先更新mask特征、再影响后续查询更新的插入位置相符。

反馈开启相对关闭在五个长度上的CAL总体AP差分别为
+0.7418, +0.4224, +0.2860, +1.3455, +1.2032个百分点。
这些汇总保留所有原始输入；其中两组存在反馈之前的运行差异，下文单独核验。
在满足严格上游对齐的28组中，当前帧新增和丢失的高IoU对象观测均为9次；
按全部前缀计算分别为5次和6次。
同一物体可能在多个前缀被重复观察，这些计数不是独立物体数。
高IoU诊断按同类候选的最佳IoU是否达到0.5计数，不是官方AP或一对一身份关联。
因此不能把候选变化只叙述为新增成功案例。

阶段采样诊断使用已训练E1的六个五复访输入，在实际层级特征上复现既定采样规则。
72次调用均触发原采样限额，共360个可用阶段位置；均匀和分层规则均未漏掉可用阶段。
分层规则使各阶段采样占比的最大最小差为0，均匀规则的调用均值为0.0612。
这支持配额更均衡，但未在该面板上证明消除了阶段遗漏。
训练轨迹中E2的短时退化仍须保留，不能由采样配额性质替代任务性能。

嵌套前缀记录包含目标不可见、命中、丢失、恢复、类别改变和原始查询槽变化。
其中不可见指目标未出现在当前官方真值中，尚未区分物理缺席与标签筛选；
相邻复访间的原始查询槽改变也不等同于持久跟踪ID切换。
